% Source: https://artofproblemsolving.com/community/c3h2865360_math_problem__games__other
% Stable identity: aops-topic-2865360-q01
Suppose the point of tangency moves through an angle $\theta$ around the clock. The arc of the clock face rolled over has length $20\theta$, so the disk turns through $\dfrac{20\theta}{10}=2\theta$ relative to the line joining the two centres. That line itself has turned through $\theta$, so the disk has turned through a total of $2\theta+\theta=3\theta$ in the plane.

The arrow next points up when $3\theta=360^\circ$, i.e. $\theta=120^\circ$. Moving $120^\circ$ clockwise from $12$ o'clock reaches $4$ o'clock.
\begin{asy}
usepackage("amsmath");
size(5cm);
defaultpen(fontsize(10pt));
import graph;
draw(circle((0,0),2));dot((0,0));label("$O$",(0,0),SW);
pair P0=(0,3),P1=3*dir(-30);
draw(circle(P0,1),gray);draw(P0--P0+(0,0.8),gray,Arrow);
draw(circle(P1,1));draw(P1--P1+(0,0.8),Arrow);
draw(arc((0,0),4.3,80,-12,CW),dashed,Arrow);
draw((0,0)--P1,dotted);draw((0,0)--P0,dotted);
label("$12$",(0,1.75));label("$4$",1.75*dir(-30));
label("$120^\circ$",(0,0),5*dir(30));
label("start",P0+(1,0),E,gray);label("end",P1-(0,1),S);
\end{asy}

The answer is $\boxed{\textbf{(C)}\ \text{4 o'clock}}$.
